\chapter{Erythropoietin (EPO)}

\section*{What Is This Test?}
The erythropoietin (EPO) test measures the serum concentration of erythropoietin, the glycoprotein hormone that is the primary regulator of red blood cell production. EPO is produced chiefly by peritubular interstitial fibroblasts of the kidney in response to tissue hypoxia and acts on erythroid progenitor cells in the bone marrow, stimulating their survival, proliferation, and differentiation into mature red blood cells. A small amount is produced by hepatocytes and is the dominant source before birth and in some anephric adults \cite{brugnara2015erythropoietin}. In health, EPO secretion follows a feedback loop: anemia or hypoxia stimulates EPO release, which drives erythropoiesis and corrects the oxygen deficit. The EPO level is measured primarily to determine why the bone marrow is not responding appropriately to an anemia (to distinguish anemias caused by EPO deficiency, such as chronic kidney disease, from those with an appropriate or elevated EPO response, such as hemolysis) and to distinguish primary polycythemia (polycythemia vera), in which EPO is inappropriately low, from secondary polycythemia, in which EPO is appropriately elevated. Because the hormone's expected concentration is inversely related to the hemoglobin level, the test is always interpreted by comparing the measured EPO to the expected EPO for the patient's degree of anemia.

\section*{Alternative Names}
Serum Erythropoietin; Erythropoietin Level; EPO Assay; Erythropoiesis-Stimulating Hormone; EP; Erythropoietin Hormone Test

\section*{Principle}
Serum EPO is measured by automated immunoassay, most commonly enzyme-linked immunosorbent assay (ELISA) or chemiluminescent immunoassay (CLIA), using monoclonal or polyclonal antibodies directed against recombinant human EPO (epoetin). The assay detects the intact circulating hormone, which is heavily glycosylated and exists as a heterogeneous mixture of isoforms with varying degrees of sialylation. Results are reported in mIU/mL, calibrated against the World Health Organization International Reference Preparation for EPO \cite{tietz2023textbook}. Because EPO production is regulated by oxygen sensing through the hypoxia-inducible factor (HIF) pathway, the physiological interpretation rests on the inverse relationship between EPO and hemoglobin: at a given degree of anemia there is an expected EPO concentration, and the measured value is judged as appropriate, inappropriately low, or inappropriately high relative to that expectation. The reference relationship is approximately log-linear; as the hematocrit falls, the expected EPO rises logarithmically. In secondary polycythemia (hypoxia-driven), EPO is elevated; in polycythemia vera (a JAK2-mutation-driven myeloproliferative neoplasm), EPO is characteristically low or inappropriately normal despite erythrocytosis \cite{kralovics2005gain}.

\section*{History}
The existence of a humoral factor controlling red blood cell production was hypothesized in 1906 by Paul Carnot, who proposed that anemia released a factor stimulating blood formation. In 1957, Leon Jacobson demonstrated that the kidney was the major site of EPO production by showing that nephrectomized animals failed to mount an erythropoietin response to hypoxia. Erythropoietin was purified in 1977 by Eugene Goldwasser and Tsuneo Miyake from the urine of patients with aplastic anemia, a feat requiring 2,550 liters of urine \cite{miyake1977purification}. The human EPO gene was cloned in 1985 by Fu-Kuen Lin at Amgen and independently by groups at Genetics Institute, enabling recombinant production \cite{lin1985cloning}. Recombinant human erythropoietin (epoetin alfa) was approved by the FDA in 1989 and revolutionized the treatment of anemia of chronic kidney disease and chemotherapy-associated anemia. The discovery that the erythropoietin receptor signals through JAK2 led to the 2005 finding that the somatic JAK2 V617F mutation underlies most cases of polycythemia vera, cementing the role of the EPO level in differentiating primary from secondary erythrocytosis. The Nobel Prize in Physiology or Medicine 2019 (Kaelin, Ratcliffe, Semenza) honored the discovery of the oxygen-sensing HIF pathway that regulates EPO production.

\section*{Why This Test Is Ordered}
\begin{itemize}
  \item \textbf{Evaluation of anemia when the cause is not apparent:} To determine whether erythropoietin production is appropriate for the degree of anemia, helping distinguish EPO-deficiency states (chronic kidney disease, anemia of chronic disease) from bone marrow or nutritional causes
  \item \textbf{Evaluation of polycythemia (elevated hemoglobin/hematocrit):} An inappropriately low EPO suggests polycythemia vera; an appropriately elevated EPO suggests secondary polycythemia (hypoxia, high altitude, smoking, chronic lung disease, cardiac shunts, high-affinity hemoglobins, and rare tumors secreting EPO)
  \item \textbf{Assessment of chronic kidney disease:} To document EPO deficiency when anemia is out of proportion to the degree of renal impairment \cite{kdoqi2007kdoqi}
  \item \textbf{Evaluation of the anemia of chronic disease:} Reduced EPO response to anemia due to inflammatory cytokine (interleukin-6, hepcidin) effects
  \item \textbf{Investigation of unexplained anemia in patients with normal kidney function:} Including suspected bone marrow failure, nutritional anemias, and hemolysis, where EPO should be appropriately elevated
  \item \textbf{Monitoring before starting erythropoiesis-stimulating agent therapy:} Baseline assessment in selected patients
\end{itemize}

\section*{Primary vs Secondary Test}
The EPO level is a secondary (reflex) test. It is ordered when anemia or polycythemia is identified on a CBC, to determine the cause of the marrow's inappropriate response. It is always interpreted together with the hemoglobin/hematocrit, reticulocyte count, iron studies, and, in polycythemia, the JAK2 mutation and oxygen saturation.

\section*{How This Test Is Performed}
A venous blood sample is collected in a red-top (serum separator, SST) tube, typically 3--5 mL. The sample may be drawn at any time of day; EPO has a circadian variation with a morning peak, but this is small relative to the clinical interpretation range. The serum is analyzed by immunoassay, with results available within 1--3 days (send-out assays are common) \cite{fischbach2023manual}. Because the interpretation depends on the simultaneous hemoglobin/hematocrit, a CBC should be drawn at the same time. In polycythemia evaluation, oxygen saturation (pulse oximetry or arterial blood gas) is measured concurrently, and the JAK2 V617F mutation test is ordered when polycythemia vera is suspected.

\section*{What Preparation Is Needed}
No fasting is required. The patient should inform the healthcare provider of: any history of blood transfusion, recent significant blood loss, iron deficiency or iron therapy, kidney disease, altitude exposure, smoking, chronic lung disease or congenital heart disease, and use of erythropoiesis-stimulating agents (epoetin alfa, darbepoetin) or testosterone (which stimulates erythropoiesis). Recent transfusion or acute bleeding can transiently suppress or elevate EPO and should be disclosed.

\section*{Contraindications and Precautions}
There are no absolute contraindications to the blood draw. Important interpretive cautions: (1) EPO must always be interpreted against the hemoglobin/hematocrit --- a ``normal'' EPO level in a severely anemic patient is actually inappropriately low; (2) recent blood transfusion raises the hematocrit and suppresses endogenous EPO, transiently lowering the measured level; (3) EPO is elevated at high altitude and in smokers with carboxyhemoglobinemia; (4) hemolyzed specimens may give falsely low results in some immunoassays; (5) athletic doping with exogenous EPO or blood transfusion suppresses endogenous EPO to very low levels; (6) pregnancy elevates EPO; (7) the distinction between polycythemia vera and secondary polycythemia also relies on JAK2 testing, oxygen saturation, and spleen size (splenomegaly supports PV), so a single EPO level is not diagnostic.

\section*{Risks}
Risks are those of routine venipuncture: mild pain, bruising, small hematoma at the collection site, lightheadedness or vasovagal syncope, and extremely rarely, infection or nerve injury.

\section*{What Happens After the Test}
The result is interpreted by comparing the measured EPO to the expected EPO for the patient's hemoglobin/hematocrit \cite{kaushansky2021williams}. In an anemic patient, an inappropriately low EPO points toward chronic kidney disease, anemia of chronic disease, or endocrinopathy, and guides therapy toward correction of the underlying condition or, in CKD, use of an erythropoiesis-stimulating agent after iron stores are replete. An appropriately elevated EPO with anemia supports hemolysis, acute blood loss, or bone marrow stress, prompting a hemolysis and iron workup. In a patient with erythrocytosis, a low EPO leads to JAK2 V617F testing and evaluation for polycythemia vera, whereas an elevated EPO leads to investigation of hypoxia (sleep apnea, chronic lung disease, cardiac shunt), hemoglobin variants with high oxygen affinity, or rare EPO-secreting tumors. Follow-up typically includes serial CBCs, reticulocyte count, iron studies, and oxygen saturation.

\section*{Understanding Results}

\subsection*{Below Normal (inappropriately low for the degree of anemia or erythrocytosis)}
\begin{itemize}
  \item \textbf{Chronic kidney disease:} The classic EPO-deficiency anemia; EPO production falls as nephron mass declines
  \item \textbf{Polycythemia vera:} EPO is characteristically low or inappropriately normal despite erythrocytosis because erythropoiesis is driven by the JAK2 mutation rather than by EPO
  \item \textbf{Anemia of chronic disease:} Inflammatory cytokines (interleukin-6) suppress EPO production and action; EPO is lower than expected for the hemoglobin
  \item \textbf{Exogenous EPO or blood doping:} Suppressed endogenous EPO
  \item \textbf{Hypothyroidism:} Reduced metabolic demand and EPO production
\end{itemize}

\subsection*{Above Normal (appropriately or inappropriately elevated)}
\begin{itemize}
  \item \textbf{Hemolytic anemia:} EPO rises appropriately in response to anemia
  \item \textbf{Acute blood loss or recent recovery:} EPO surges within hours of hemorrhage
  \item \textbf{Secondary polycythemia:} Tissue hypoxia (high altitude, smoking/carbon monoxide, chronic obstructive pulmonary disease, sleep apnea, cyanotic congenital heart disease, high-affinity hemoglobin variants) drives EPO elevation
  \item \textbf{Renal pathology with local hypoxia:} Renal cysts, renal cell carcinoma, and other tumors may secrete EPO ectopically
  \item \textbf{Pregnancy and high altitude:} Physiological EPO elevation
  \item \textbf{Iron deficiency without anemia (early):} Mildly elevated EPO as erythropoiesis becomes less efficient
\end{itemize}

\section*{Normal Results}
\begin{itemize}
  \item \textbf{Normal range:} Approximately 4--25 mIU/mL (laboratory-dependent)
  \item \textbf{Expected relationship:} EPO rises logarithmically as hemoglobin falls below normal; at a hemoglobin of approximately 10 g/dL the expected EPO is roughly 30--100 mIU/mL, and values of several hundred to >1,000 mIU/mL are expected in severe anemia (e.g., hemoglobin $<$8 g/dL) from hemolysis or blood loss
  \item \textbf{Polycythemia vera:} EPO below the reference range or inappropriately normal for the elevated hematocrit
  \item \textbf{Secondary polycythemia:} EPO elevated or high-normal
\end{itemize}

\section*{Follow-up Tests}
\begin{itemize}
  \item \textbf{Complete blood count with reticulocyte count:} To interpret EPO against the degree of anemia and marrow response
  \item \textbf{Iron studies (ferritin, serum iron, TIBC, transferrin saturation):} Iron deficiency blunts the response to EPO and is required before and during ESA therapy
  \item \textbf{Serum creatinine and estimated GFR:} To assess renal function and EPO-deficiency anemia
  \item \textbf{Vitamin B12, folate, and methylmalonic acid:} To exclude nutritional anemias when EPO response is low
  \item \textbf{JAK2 V617F mutation:} When polycythemia vera is suspected (low EPO with erythrocytosis)
  \item \textbf{Arterial blood gas or pulse oximetry:} To evaluate hypoxia as a cause of secondary polycythemia
  \item \textbf{Sleep study:} Obstructive sleep apnea is a common cause of secondary polycythemia
  \item \textbf{Hemoglobin electrophoresis and oxygen affinity studies:} For high-affinity hemoglobin variants causing familial erythrocytosis
  \item \textbf{Abdominal imaging:} When ectopic EPO secretion from renal or hepatic tumors is suspected
  \item \textbf{Bone marrow aspiration and biopsy:} In myeloproliferative neoplasms and unexplained marrow failure
\end{itemize}

\section*{References}
\bibliographystyle{unsrtnat}
\bibliography{references}

\clearpage
